% Document setup
\RequirePackage{fix-cm}
\documentclass{article}
\usepackage{preamble}

%===============================%
%=============Inputs============%
%===============================%
\newcommand{\studentname}{Aaditya Shrivastava}

%===============================%
%====Make Solutions Visible=====%
%===============================%
\showsolutionstrue

%===============================%
%=============Figures===========%
%===============================%
% Remember to start figures with 

% \begin{figure}[H]

% the [H] is critical!


%==============================%
%=====Beginning of Document====%
%==============================%
\begin{document}
\pagestyle{mypagestyle}




%%%%%%%%%%%%%%%%%%%%%%%%%%%%%%%%%%%%%%%%%%%%%%%%%%%%%%%%%%%%%%%%%%%%%%%%%%%%%%%%%%%%%%%%%%%%%%%%%%%%%%%%%%%%%%%%%%%%%%%%%%
% Project 7: Aerial Firefighting
%%%%%%%%%%%%%%%%%%%%%%%%%%%%%%%%%%%%%%%%%%%%%%%%%%%%%%%%%%%%%%%%%%%%%%%%%%%%%%%%%%%%%%%%%%%%%%%%%%%%%%%%%%%%%%%%%%%%%%%%%%

\project{7}{Aerial Firefighting (100 points)}{Wed. Apr. 22 @ 11:59PM PST}


\subsection*{Motivation}
Aerial firefighting is a critical tool in modern wildfire suppression, particularly as climate change drives longer fire seasons and larger burn areas. Fixed-wing aircraft and air tankers can deliver fire retardant over steep, inaccessible terrain far faster than ground crews, but their effectiveness depends heavily on \textit{when}, \textit{where}, and \textit{how} the retardant is released. Poor drop geometry---wrong altitude, speed, or heading---can scatter retardant outside the fire perimeter entirely. Optimizing these parameters manually is infeasible given the high-dimensional design space and the complex, coupled physics involved. This project uses physics-based simulation paired with a genetic algorithm to systematically search for drop configurations that maximize retardant coverage over a target fire zone.

\subsection*{Governing Equations}
The simulation couples aircraft path kinematics with Lagrangian particle dynamics. The aircraft follows a parametric sine-wave trajectory $\br(u) = (u + x_0,\; y_0,\; A\sin(2\pi u/\lambda) + z_0)^T$, advanced each timestep by $\Delta u = V\Delta t / \|\br'(u)\|$ to enforce constant airspeed. Retardant droplets are ejected from four nozzles with velocity $\bv_\text{drop} = \bv_\text{aircraft} + v_e\,\hat{\bn}_\text{spray}$, where the spray direction $\hat{\bn}_\text{spray}$ is stochastically perturbed to model nozzle spread. Once airborne, each droplet obeys Newton's second law,
\begin{equation}
    m\ddot{\br} = \bff_\text{drag} + m\bg,
\end{equation}
where the drag force is $\bff_\text{drag} = \frac{1}{2}\rho_\text{air} C_D (\pi r_p^2)\|\bv_\text{rel}\|\bv_\text{rel}$ with $\bv_\text{rel} = (\bv_\text{wind} + \bv_\text{up}) - \bv_\text{drop}$. Above the fire footprint, a thermal updraft $v_\text{up}(y) = v_{\text{up,max}}\exp(-a_{\text{decay}}\max(0,y))$ opposes descent. $v_{\text{up,max}}$ and $a_{\text{decay}}$ are the maximum velocity of the updraft (closest to the fire) and the decay rate as the air moves away from the fire, respectively. The equations of motion are integrated via forward Euler, and a drop is considered successful if it lands within the prescribed fire boundary. The optimization objective is to minimize the cost 
\begin{equation}
    \Pi = W_1\left(1 - \frac{N_\text{hit}}{N_\text{total}} \right) + W_2\left(1 - \frac{y_0}{y_{\text{max}}} \right)
\end{equation}
which maximizes the fraction of particles that hit the fire region, $N_{\text{hit}}$ while flying as high as possible, $y_0$. Note, starting height of the plane is constant in this project.



\subsection*{Modifications and Comments on Original Paper}
Several modifications were made for this assignment, but you can read the full paper that inspired this project \href{https://msol.berkeley.edu/wp-content/uploads/2023/10/181.pdf}{here}. First, the flight path of the plane is defined by a parametric curve, not calculated by variables such as radial velocity. We feel a curve is more intuitive and can be clearly communicated to a pilot when they fly the mission. The domain of fire is now a continuous rectangle, rather than a region of discrete targets. This was for simplicity of the simulation. This led to a modified cost function, since the original maximized particles in the fire region and targets hit. Now the objective is to maximize particles in the region and the height of the aircraft. Flying higher is desirable to enhance pilot safety. 


\subsection*{Deliverables}
Fill in the answer boxes on this report. Once your 
\begin{enumerate}

%%%%%%%%%%%%%%%%%%%%%%%%%%%%%%%
% Problem 1
%%%%%%%%%%%%%%%%%%%%%%%%%%%%%%%

\item[1.] \textbf{Particle physics -- drag and updraft (15 points):}
In flight, each active droplet obeys
\begin{equation}
    m\,\ddot{\br} = \bff_\text{drag} + m\,\bg,
    \label{eq:eom}
\end{equation}
where the aerodynamic drag force is
\begin{equation}
    \bff_\text{drag} = \tfrac{1}{2}\,\rho_\text{air}\,C_D\,(\pi r^2)\,\|\bv_\text{rel}\|\,\bv_\text{rel},
    \quad \bv_\text{rel} = (\bv_\text{wind} + \bv_\text{up}) - \bv_\text{drop}.
    \label{eq:drag}
\end{equation}
In the original paper, the fire-induced updraft velocity (applied only above the fire footprint) is
\begin{equation}
    v_\text{up}(y) = 30\,\exp\!\left(-0.005\,\max(0,\,y)\right)\,\hat{\be}_y,
    \label{eq:updraft}
\end{equation}
where $y$ is the droplet height above the ground.

\begin{enumerate}[(i)]
    \item The simulation uses forward Euler integration. Write out the discrete update equations for $\bv$ and $\br$ at timestep $t+1$ from timestep $t$.
    \item What physical phenomenon does \cref{eq:updraft} model? How does the exponential decay in $y$ reflect the behavior of a real fire plume?
    \item Identify one limitation of the drag model in \cref{eq:drag} (e.g.\ regarding the droplet Reynolds number regime or shape assumption).
\end{enumerate}

\begin{answer}
% Problem 1 answers

\textbf{(i) Forward Euler discrete update equations:}

Given acceleration $\ba^t = \bff_\text{drag}^t / m + \bg$, the forward Euler updates are:
\begin{align*}
    \bv^{t+1} &= \bv^t + \ba^t \,\Delta t \\
    \br^{t+1} &= \br^t + \bv^{t+1} \,\Delta t
\end{align*}
where $\ba^t = \bff_\text{drag}(\bv_\text{rel}^t)/(\rho_\text{ret}\,v_p) + \bg$ and $\bv_\text{rel}^t = (\bv_\text{wind} + \bv_\text{up}) - \bv^t$.

\textbf{(ii) Physical phenomenon modeled by Equation~(5):}

Equation~(5) models the fire-induced thermal updraft: hot combustion gases rise above the fire and entrain surrounding air into a buoyant plume. The exponential decay in height $y$ reflects the fact that the vertical momentum of the plume dissipates as it rises, mixing with cooler ambient air and losing buoyancy; the plume is strongest at ground level and weakens with altitude.

\textbf{(iii) One limitation of the drag model in Equation~(4):}

The model assumes a perfectly spherical droplet with a fixed projected area $\pi r_p^2$ and uses a scalar drag coefficient $C_D$ that depends only on Reynolds number. In reality, retardant droplets deform and may break up at high relative speeds (droplet Weber number effects), causing the effective drag area and $C_D$ to differ substantially from the spherical assumption.
\end{answer}


%%%%%%%%%%%%%%%%%%%%%%%%%%%%%%%
% Problem 2
%%%%%%%%%%%%%%%%%%%%%%%%%%%%%%%

\item[2. (a)] \textbf{Design vector $\BLambda$ and parameter bounds (5 points):}
The design vector
\begin{equation}
    \BLambda = \bigl[V,\; x_0,\; y_0,\; z_0,\; r_p,\; A_s,\; t_\text{start},\; t_\text{end},\; \dot{Q},\; v_e,\; A_c,\; \lambda_c\bigr]
    \label{eq:lam}
\end{equation}
has 12 components.
The parameter ranges (from \texttt{geneticalgorithm.py}) are summarized in \cref{tab:lam}.

\begin{table}[H]
\centering
\caption{Design vector components and their admissible ranges.}
\label{tab:lam}
\begin{tabular}{clccc}
\toprule
Index & Symbol & Description & Lower Bound & Upper Bound \\
\midrule
0  & $V$             & Aircraft speed              & $50$ m/s & $100$ m/s \\
1  & $x_0$           & Initial $x$-position        & $0$ m              & $0$ m (fixed)       \\
2  & $y_0$           & Initial $y$-position        & $50$ m              & $250$ m             \\
3  & $z_0$           & Initial $z$-position        & $-100$ m           & $100$ m             \\
4  & $r_p$           & Droplet radius              & $0.001$ m          & $0.01$ m            \\
5  & $A_s$           & Spray amplitude             & $0$                & $0.5$               \\
6  & $t_\text{start}$& Release start time          & $0$ s              & $0.25\,T$           \\
7  & $t_\text{end}$  & Release end time            & $0.25\,T$          & $T$                 \\
8  & $\dot{Q}$       & Retardant flow rate         & $2$ m$^3$/s        & $10$ m$^3$/s        \\
9  & $v_e$           & Ejection speed              & $0$ m/s            & $150$ m/s           \\
10 & $A_c$           & Curve amplitude             & $1$ m              & $100$ m             \\
11 & $\lambda_c$     & Curve wavelength            & $500$ m            & $2000$ m            \\
\bottomrule
\end{tabular}
\end{table}

In a sentence, why is $x_0$ fixed at $0$ rather than included in the optimization?

\begin{answer}
% Problem 2(a) answer
$x_0$ is fixed at 0 because the aircraft's absolute starting $x$-position does not affect drop coverage---only the relative geometry of the flight path over the fire zone matters, so fixing $x_0 = 0$ removes a redundant degree of freedom and prevents the optimizer from wasting budget on a parameter that shifts the entire trajectory without changing its shape.
\end{answer}

\item[2. (b)] \textbf{Implement \texttt{student\_simulation.py} (40 points):}
Open \texttt{student\_simulation.py}. Rename the file \texttt{simulation.py}. Follow the instructions below then run \texttt{main.py} to generate figures and numerical results.
When finished, confirm your implementation reproduces a cost of \textbf{0.1800} when run with the fixed design vector:
\begin{center}
\texttt{LAM = [69.3, 0.0, 175.1, -11.92, 0.00429, 0.183, 0.02159, 0.6921, 8.21e-6, 48.78, 48.5, 1520.7]}
\end{center}

\begin{enumerate}[(i)]
    \item \textbf{Step~1 -- \texttt{\_update\_aircraft\_state}:}
          The aircraft travels along a parametric curve at constant airspeed $V$.
          Fill in: (a) the position and tangent from the curve at the current parameter $u$,
          (b) the step size $\Delta u = V\,\Delta t\,/\,\|\bt\|$ so the aircraft advances at constant speed, and
          (c) the new \texttt{airplane.pos}.

    \item \textbf{Step~2 -- \texttt{\_compute\_nozzle\_positions}:}
          Four nozzles are arranged around the aircraft.
          Fill in the lateral and aft offsets for nozzles 1, 2, and 3, and the
          vertical drop from the sprayer base to its end.
          Reference the nozzle layout comments in the file.

    \item \textbf{Step~3 -- \texttt{\_compute\_spray\_velocities}:}
          The spray direction is the nozzle axis perturbed by random noise.
          Fill in: (a) the perturbed spray direction
          $\hat{\bn}_\text{base} + A_s\,\hat{\bn}_\text{rand}$, and
          (b) the final droplet velocity $\bv_\text{nozzle} + v_e\,\hat{\bn}_\text{spray}$.

    \item \textbf{Step~4 -- \texttt{\_release\_droplets}:}
          Fill in the one-timestep position pre-integration
          (\texttt{sprayer\_end + DT * velocities}) so each droplet starts
          one step ahead of the nozzle.

    \item \textbf{Step~5 -- \texttt{\_in\_fire\_zone}:}
          Fill in the squared distance from each particle to the circle center
          $(x - x_c)^2 + (z - z_c)^2$, and the comparison against $R^2$.
          The box case is already provided for reference.

    \item \textbf{Step~6 -- \texttt{\_compute\_updraft}:}
          Fill in the updraft speed for particles inside the fire zone:
          \begin{equation*}
              v_\text{up}(y) = V_\text{up}\,\exp\!\bigl(-\kappa\,\max(0,\,y)\bigr),
          \end{equation*}
          where $V_\text{up} = $ \texttt{UPDRAFT\_VELOCITY} and $\kappa = $ \texttt{UPDRAFT\_DECAY}.

    \item \textbf{Step~7 -- \texttt{\_compute\_drag}:}
          Fill in the Reynolds number
          $\mathrm{Re} = 2\,r_p\,\rho_a\,|\bv_\text{rel}|\,/\,\mu$.
          The Chow piecewise $C_D$ table and the drag force assembly are already provided.

    \item \textbf{Step~8 -- \texttt{\_step\_physics}:}
          Fill in three lines using Newton's second law and Semi-Implicit Euler:
          \begin{align*}
              \bv_\text{rel} &= (\bv_\text{wind} + \bv_\text{up}) - \bv_\text{curr}, \\
              \ba             &= \bff_\text{drag}\,/\,(\rho_\text{ret}\,v_p) + \bg, \\
              \bv             &\mathrel{+}= \ba\,\Delta t, \\ 
              \br              &\mathrel{+}= \bv\,\Delta t.
          \end{align*}

    \item \textbf{Step~9 -- \texttt{simulate\_path\_with\_nozzles}:}
          Fill in the physics mask passed to \texttt{\_step\_physics}:
          particles that are \texttt{active} and \textbf{not} yet \texttt{landed}.

    \item \textbf{Step~10 -- \texttt{calculate\_cost}:}
          Fill in the cost function below:
          \begin{equation*} \label{eq:cost}
              \Pi = W_1\!\left(1 - \frac{N_\text{hit}}{N_\text{total}}\right)
                  + W_2\!\left(1 - \frac{y_0}{y_{\max}}\right).
          \end{equation*}
          where $N_\text{hit}$ is the number of particles that land within the fire boundary,  $N_\text{total}$ is the total number of particles released, $y_0$ is the fixed height of the aircraft, and $y_{\text{max}}$ is the maximum possible height we allow the plane to fly.
To help with debugging, confirm the cost printed is 0.1800 and uncomment the animation.
\end{enumerate}



\item[2. (c)] \textbf{GA convergence (10 points):}
Run the genetic algorithm (\texttt{geneticalgorithm.py}) using your completed simulation and plot the best cost,
average cost, and parents' average cost versus generation number.
Report the best design vector $\BLambda^*$ found and its corresponding cost $\Pi(\BLambda^*)$.
Include a screenshot of the animation for the optimized design.

\begin{answer}

\begin{figure}[H]
    \centering
    \includegraphics[width=0.90\textwidth]{ga_convergence.png}
    \caption{GA convergence over 20 generations. Blue: best cost in population; green: parents' average cost (top $P=6$); orange: full population average cost ($S=20$).}
    \label{fig:ga_convergence}
\end{figure}

All three metrics begin near 0.87--0.93 (random initial population), drop sharply around generation 7--8 as the GA discovers a good region of the design space, then gradually converge by generation 20. The persistent gap between the orange (average) and green (parents avg) curves shows that random immigrants continually introduce high-cost designs, keeping the full-population average elevated while the best solution improves.

\textbf{Best design vector:}
\begin{equation*}
    \BLambda^* = [63.18,\ 0,\ 177.23,\ 11.88,\ 0.00354,\ 0.096,\ 0.031,\ 0.681,\ 8.65,\ 91.12,\ 4.88,\ 1558.0]
\end{equation*}

\begin{table}[H]
\centering
\begin{tabular}{clcc}
\toprule
Index & Symbol & Value & Unit \\
\midrule
0  & $V$              & 63.18  & m/s \\
1  & $x_0$            & 0      & m (fixed) \\
2  & $y_0$            & 177.23 & m \\
3  & $z_0$            & 11.88  & m \\
4  & $r_p$            & 0.00354 & m \\
5  & $A_s$            & 0.096  & -- \\
6  & $t_\text{start}$ & 0.031  & fraction of $T$ \\
7  & $t_\text{end}$   & 0.681  & fraction of $T$ \\
8  & $\dot{Q}$        & 8.65   & m$^3$/s \\
9  & $v_e$            & 91.12  & m/s \\
10 & $A_c$            & 4.88   & m \\
11 & $\lambda_c$      & 1558.0 & m \\
\bottomrule
\end{tabular}
\caption{Optimal design vector $\BLambda^*$ components.}
\end{table}

\textbf{Optimal cost:} $\Pi(\BLambda^*) = \mathbf{0.1310}$

Breaking this down: $W_2(1 - y_0/y_\text{max}) = 0.1(1 - 177.23/250) = 0.029$, so the coverage term contributes $0.1310 - 0.029 = 0.102$, giving $N_\text{hit}/N_\text{total} = 1 - 0.102/0.9 \approx 88.7\%$ of tracked droplets landing inside the fire zone.

\textbf{Animation screenshot of optimized design:}

\begin{figure}[H]
    \centering
    \includegraphics[width=0.72\textwidth]{ga_result_screenshot.png}
    \caption{Final frame of the firefighting simulation with $\BLambda^*$. Blue dots mark retardant landing positions; the red rectangle is the fire zone ($x \in [200,300]$\,m, $z \in [-50,50]$\,m). The majority of droplets land inside the fire boundary.}
    \label{fig:animation}
\end{figure}

\end{answer}


%%%%%%%%%%%%%%%%%%%%%%%%%%%%%%%
% Problem 3
%%%%%%%%%%%%%%%%%%%%%%%%%%%%%%%

\item[3.] \textbf{Particle budget and flow rate (5 points):}
The total number of particles in the simulation is
\begin{equation}
    N_\text{total} = \left\lfloor \frac{V_\text{tank}}{v_p} \right\rfloor, \qquad v_p = \tfrac{4}{3}\pi r_p^3,
    \label{eq:ntotal}
\end{equation}
and the number released per timestep is $\lfloor \dot{Q}\,\Delta t / v_p \rfloor$.


The simulation has been updated to use more realistic droplet radii ($0.001$--$0.01$ m, or $1$--$10$ mm diameter) compared to earlier versions that used much larger droplets. To keep particle counts computationally manageable with these realistic smaller droplets, each of the smaller tracked droplets is a "super-particle" representing ($10^6$) droplets. 

Why is using realistic droplet sizes with a full-size tank problematic computationally, and what trade-off does the super-particle method introduce?

\begin{answer}
% Problem 3 answer
With a 10\,m$^3$ tank and droplet radii of $r_p = 0.001$--$0.01$\,m (volumes $v_p \sim 4\times10^{-9}$--$4\times10^{-6}$\,m$^3$), the total particle count $N_\text{total} = \lfloor V_\text{tank}/v_p \rfloor$ reaches $\sim\!2.5\times10^{9}$ for the smallest droplets---far too many to integrate individually each timestep without prohibitive memory and CPU cost. The super-particle method collapses $10^6$ physical droplets into one tracked particle, reducing $N_\text{total}$ to a computationally feasible $\mathcal{O}(10^3)$; the trade-off is that it sacrifices statistical resolution: rare events (e.g.\ a few droplets landing precisely on the fire boundary) are averaged away, and the discrete hit count $N_\text{hit}$ becomes a coarser estimate of actual retardant coverage.
\end{answer}

%%%%%%%%%%%%%%%%%%%%%%%%%%%%%%%
% Problem 4
%%%%%%%%%%%%%%%%%%%%%%%%%%%%%%%
\item[4.] \textbf{Mini-Sensitivity Analysis (25 Points):}
A sensitivity analysis tells us how much an outcome changes with respect to changes in its input variable(s). Mathworks offers a nice summary \href{https://www.mathworks.com/help/sldo/ug/what-is-sensitivity-analysis.html}{here}. In this problem we examine how ``sensitive" our cost is to the drop velocity of the particles and the drop velocity of the plane. We have already coded the sample. First, run the code with 
\begin{lstlisting}
    ga_results_path = ''
\end{lstlisting}
and add the plots below. Then, fill in the path to your pickle of your simulation's genetic algorithm results (should just be ga\_results.pkl). Add the plots below, clearly label each set of plots.\\
Looking at the output plots, argue why a genetic algorithm is a good choice of optimizer for this problem.

\begin{answer}

\textbf{Set 1 — Default baseline LAM} ($\dot{Q} = 8.21\times10^{-6}$\,m$^3$/s, $v_e = 48.78$\,m/s):

\begin{figure}[H]
    \centering
    \begin{subfigure}[b]{0.48\textwidth}
        \includegraphics[width=\textwidth]{sensitivity_surface_default.png}
        \caption{3D surface — default LAM}
    \end{subfigure}
    \hfill
    \begin{subfigure}[b]{0.48\textwidth}
        \includegraphics[width=\textwidth]{sensitivity_heatmap_default.png}
        \caption{Heatmap — default LAM}
    \end{subfigure}
    \caption{Sensitivity sweep centered on the default baseline LAM. Cost is nearly flat ($\approx 0.90$--$0.96$) across the entire range of plane velocity and ejection speed.}
    \label{fig:sens_default}
\end{figure}

\textbf{Set 2 — GA-optimized LAM} ($\BLambda^*$, $v_e = 91.12$\,m/s, $V = 63.18$\,m/s):

\begin{figure}[H]
    \centering
    \begin{subfigure}[b]{0.48\textwidth}
        \includegraphics[width=\textwidth]{sensitivity_surface_ga.png}
        \caption{3D surface — GA best LAM}
    \end{subfigure}
    \hfill
    \begin{subfigure}[b]{0.48\textwidth}
        \includegraphics[width=\textwidth]{sensitivity_heatmap_ga.png}
        \caption{Heatmap — GA best LAM}
    \end{subfigure}
    \caption{Sensitivity sweep centered on the GA-optimized $\BLambda^*$. The cost landscape is highly nonlinear with a deep minimum (cost $\approx 0.025$) near $V \approx 55$--$65$\,m/s and $v_e \approx 0$--$30$\,m/s, surrounded by steep ridges and multiple local minima.}
    \label{fig:sens_ga}
\end{figure}

\textbf{Why a genetic algorithm is the right optimizer:}

The two sensitivity sweeps together make the case for the GA on three grounds:

\begin{enumerate}
    \item \textbf{Multimodal, non-convex landscape.} The GA-optimized surface (Figure~\ref{fig:sens_ga}) contains multiple distinct valleys separated by ridges of high cost. A gradient-based optimizer (e.g., gradient descent, Newton's method) follows the local slope and converges to whichever local minimum it starts near, with no mechanism to escape. The GA avoids this by maintaining a \emph{population} of $S=20$ diverse designs simultaneously exploring different regions, plus random immigrants every generation that inject fresh solutions into unexplored areas.

    \item \textbf{Non-smooth, discontinuous cost.} The cost function is driven by $N_\text{hit}$, which counts whether each droplet lands \emph{inside or outside} a hard boundary — a binary outcome. This produces a nearly piecewise-constant surface with sharp edges where the gradient is either zero or undefined. Gradient-based methods require a smooth, differentiable objective; they cannot function here. The GA requires only cost \emph{evaluations}, not derivatives.

    \item \textbf{High-dimensional coupled design space.} The 12-parameter $\BLambda$ exhibits strong nonlinear interactions: ejection speed interacts with droplet radius through drag, which depends on altitude, timing, and wind. The flat default-LAM surface (Figure~\ref{fig:sens_default}) demonstrates that most of the parameter space produces uniformly poor solutions — the good region is a narrow, hard-to-find valley. The GA's crossover operator (children = weighted average of two parents) naturally exploits these interactions by inheriting correlated parameter combinations from successful parents, whereas coordinate-wise methods adjust one parameter at a time and are blind to such correlations.
\end{enumerate}

\end{answer}

\end{enumerate}

\subsection*{Submission}
\begin{enumerate}
    \item Submit the PDF copy of your typed report on Gradescope before the end of day (11:59 pm PST).
\end{enumerate}


% End document
\end{document}